This equality is the sign of a structural failure. With four excerpts retrieved
and four documents in the corpus, recall@4 ought to be close to
\SI{100}{\percent}. That it equals accuracy@1 means the four excerpts
\emph{always come from the same document}: retrieval systematically collapses
onto a single source.

\subsection{Diagnosis: a cross-lingual failure}

Examining the corpus provides the explanation. It is \textbf{bilingual}, and the
split is clear-cut: the two files that actually document OFFBEAT — the structure
of a case and the list of commands — are in \textbf{English}, since they come
from the code repository; the documents added by the project are in
\textbf{French}.

Yet the agent is driven in French. The embedding model in use then matches
French questions to French documents, irrespective of their content. The OFFBEAT
documentation proper is never reached.

The control experiment removes all doubt: the same questions asked in English
correctly retrieve the English documents.

\begin{table}[htbp]
\centering
\caption[Effect of language on retrieval]{Effect of the question language on
documentary retrieval, with corpus and model unchanged. The retrieval mechanism
works; it is crossing the language barrier that fails.}
\label{tab:rag-langue}
\begin{tabularx}{0.92\textwidth}{@{}Xcc@{}}
\toprule
\textbf{Question language} & \textbf{Accuracy@1} & \textbf{Recall@4}\\
\midrule
French $\rightarrow$ English documents & \textcolor{cAlert}{\textbf{0/11}} &
  \textcolor{cAlert}{\textbf{0/11}}\\
English $\rightarrow$ English documents & 9/11 (\SI{82}{\percent}) &
  \textbf{11/11 (\SI{100}{\percent})}\\
\bottomrule
\end{tabularx}
\end{table}

\subsection{Correction and re-measurement}

The natural fix is not to translate the documentation — it belongs to the
upstream repository and would have to be maintained at every version upgrade —
but to replace the embedding model with a \textbf{multilingual} one, which
projects a French text and its English translation to the same place in the
vector space. The corpus and the index were rebuilt identically, only the
embedding model changing.

\begin{table}[htbp]
\centering
\caption[Documentary retrieval quality]{Documentary retrieval quality before and
after switching to a multilingual embedding model. Corpus, questions and
chunking identical.}
\label{tab:rag}
\begin{tabularx}{0.92\textwidth}{@{}Xcc@{}}
\toprule
\textbf{Indicator} & \textbf{Initial model} & \textbf{Multilingual model}\\
\midrule
Accuracy@1 & \SI{27}{\percent} & \textbf{\SI{87}{\percent}}\\
Recall@4 & \SI{27}{\percent} & \textbf{\SI{100}{\percent}}\\
\midrule
Case structure & 0/6 & \textbf{6/6}\\
Commands & 0/5 & 3/5 (recall@4: 5/5)\\
Adding an error entry & 4/4 & 4/4\\
\bottomrule
\end{tabularx}
\end{table}

Recall@4 reaches \SI{100}{\percent}: the agent now always receives the right
document among its four excerpts. The multilingual model has become the
project's default.

This case completes the two preceding ones and forms a pattern with them: none
of these three failures — overlapping tool descriptions, degenerate training
domain, cross-lingual retrieval — was visible without measurement, and each was
fixed without rewriting any application code. What was missing was not
technique, but \emph{instrumentation}.

\section{Confronting the criterion with what the solver computes}
\label{sec:seuils}

In the architecture adopted, safety thresholds are simple constants written in a
JSON file. Checking these constants against the physical model actually solved
revealed two defects of the same nature, and suggested a deeper improvement.

\subsection{A criterion checking a limit the simulation does not use}

The cladding yield strength was set to \SI{350}{\mega\pascal} in the criteria
base. Yet the template configures the plasticity model with a limit of
\SI{250}{\mega\pascal}. The safety criterion was therefore evaluating the stress
computed by the solver \emph{against a limit that the very same solver does not
apply} — a discrepancy of \SI{100}{\mega\pascal}, that is \SI{40}{\percent}. On
the reference case, the reported margin thus went from \SI{32}{\percent} in
reality to \SI{23}{\percent} as displayed.

The fix does not consist in copying across the right constant, since that would
reproduce the same risk of drift. OFFBEAT writes the yield strength as a
\textbf{field}, \texttt{sigmaY}, cell by cell: the criteria base was therefore
extended so that a limit may designate a field (\texttt{limit\_field}) rather
than a constant. The ratio to the limit is then computed \emph{point by point}
and only then reduced:
\[
  \max_{x}\ \frac{|\sigma_{\theta\theta}(x)|}{\sigma_Y(x)}
\]

This point-by-point computation is not a cosmetic refinement. The stress peak and
the minimum of the limit are not at the same location; comparing the maximum of
one with the minimum of the other yields a result that corresponds to no point of
the rod. On a test case built for the purpose, the naive criterion concludes
\SI{83}{\percent} of the limit where the correct computation gives
\SI{150}{\percent} — that is, a missed exceedance.

The main benefit is architectural: the criterion now \emph{automatically}
inherits any temperature, irradiation or burnup dependence that the material
model may come to carry, without a single value having to be maintained twice.

\subsection{A criterion measuring something other than what it claimed}

The PCMI criterion states a limit on the \textbf{plastic} circumferential
strain. It was in fact reading \texttt{epsilonCyl}, the \textbf{total} strain.
The decomposition of the strains on the reference case
(table~\ref{tab:deformations}) shows the extent of the misunderstanding.

\begin{table}[htbp]
\centering
\caption[Decomposition of the cladding strain]{Decomposition of the
circumferential cladding strain at end of cycle. The PCMI criterion was reading
the total, whereas its statement bears on the plastic component alone — which is
exactly zero.}
\label{tab:deformations}
\begin{tabularx}{0.9\textwidth}{@{}Xcc@{}}
\toprule
\textbf{Component} & \textbf{Maximum value} & \textbf{Reversible?}\\
\midrule
Total (\texttt{epsilonCyl}, read by the old criterion) & \num{1.81e-3} & ---\\
\midrule
Plastic (\texttt{epsilonP}) & \textbf{exactly 0} & no\\
Creep (\texttt{epsilonCreepEq}) & \num{3.55e-3} & no\\
Elastic (\texttt{epsEl}) & \num{1.03e-3} & yes\\
Thermal (\texttt{epsTh}) & \num{2.31e-3} & yes\\
\bottomrule
\end{tabularx}
\end{table}

The criterion was therefore reporting \SI{18}{\percent} of a limit whose measured
quantity is in reality at \SI{0}{\percent}. The error erred on the conservative
side — total strain bounds plastic strain from above — but it measured the wrong
quantity, and above all it masked the essential point: the irreversible strain
actually present is the \textbf{creep} strain, almost twice the reported total,
and no criterion was monitoring it.

The PCMI criterion now reads the plastic strain, in accordance with its
statement, and a separate criterion was added for creep strain. Whether the
\SI{1}{\percent} design criterion bears on plasticity alone or on total
permanent strain remains an open question and is among the points to be settled
with the supervisor.

\subsection{Knock-on effect on the surrogate}

These corrections had a consequence that only a systematic check could reveal.
The surrogate of chapter~3 records its target quantities \emph{through} the
criteria base: making the PCMI criterion conform to its statement therefore
brought one of its three targets down to zero, plastic strain being null in
nominal operation. The target would have been degenerate, with no variation to
learn.

The surrogate now targets creep strain, and the design of experiments was
rebuilt accordingly. The outcome is better than before: creep accumulates
monotonically, where the total strain was non-monotonic, and the
cross-validation coefficient of determination rises from $0.9535$ to
$\mathbf{0.9905}$ (table~\ref{tab:surrogate}). The correct quantity turns out to
be the easier one to learn as well. The prediction of the gap closure date, the
main result of section~\ref{sec:validation-emulateur}, is unchanged.

This is one more illustration of the coupling between components believed to be
independent: a correctness fix in the threshold base silently invalidated a
learning target two layers away. It was detected only because every change was
followed by an end-to-end check.

\section{Correctness defects identified and fixed}
\label{sec:bugs}

A second form of stress-testing proved at least as fruitful: no longer
provoking crashes, but confronting obtained results with expected ones. Fourteen
defects were identified in this way, \textbf{none of which raised any visible
error} (table~\ref{tab:bugs}). These are \textit{correctness} errors, silent by
nature, and therefore particularly worrying in a safety context: the system was
producing plausible, correctly formatted, and wrong numbers.
